\documentclass[openany]{book}
\usepackage[a4paper, margin=1in]{geometry}
\usepackage[margin=0pt, noheadfoot]{geometry}
\usepackage{minted}
\setminted{breaklines, % Break lines automatically
    breakanywhere, % Allow breaks anywhere
    fontsize=\small, % Font size for the code
    frame=single, % Border style for the code block
    bgcolor=darkbackground, % Background color for the code block
    linenos, % Display line numbers
    numbersep=5pt, % Space between line numbers and code
    label=code-\thecodeblock}
\usepackage{caption}
\usepackage{tcolorbox}
\tcbuselibrary{breakable, skins, theorems, listings}
\usepackage{titlesec}
\usepackage{fancyhdr}
\usepackage{hyperref}
\usepackage{lipsum}
\usepackage{amsmath}
\usepackage{amsthm}
\usepackage{amssymb}
\usepackage{enumitem}
\usepackage{tikz}
\usepackage{tocloft}

% Colors
\usepackage{xcolor}
\definecolor{darkbackground}{HTML}{121212}
\definecolor{whitetext}{HTML}{FFFFFF}
\definecolor{cyanaccent}{HTML}{00FFFF}
\definecolor{lightcyan}{HTML}{E0FFFF}
\definecolor{lightgray}{HTML}{D3D3D3}
\definecolor{accent1}{HTML}{0047AB}
\definecolor{accent2}{HTML}{6495ED}
\definecolor{accent3}{HTML}{3F00FF}
\definecolor{myorange}{HTML}{FFA500}
\definecolor{mylightblue}{HTML}{ADD8E6}
\definecolor{mygreen}{HTML}{90EE90}

\pagecolor{darkbackground}
\color{whitetext}

\newenvironment{longlisting}{\captionsetup{type=listing}}{}

% Custom Counters
\newcounter{mydefinition}[subsection]
\renewcommand{\themydefinition}{\thesubsection.\arabic{mydefinition}}
\newcounter{myidentity}[subsection]
\renewcommand{\themyidentity}{\thesubsection.\arabic{myidentity}}
\newcounter{myproof}[subsection]
\renewcommand{\themyproof}{\thesubsection.\arabic{myproof}}
\newcounter{mynote}[subsection]
\renewcommand{\themynote}{\thesubsection.\arabic{mynote}}
\newcounter{mycodeblock}[subsection]
\renewcommand{\themycodeblock}{\thesubsection.\arabic{mycodeblock}}
\newcounter{myexercise}[subsection]
\renewcommand{\themyexercise}{\thesubsection.\arabic{myexercise}}
\newcounter{codeblock}[section]
\renewcommand{\thecodeblock}{\thesubsection.\arabic{codeblock}}


\setminted{breaklines, % Break lines automatically
    breakanywhere, % Allow breaks anywhere
    fontsize=\small, % Font size for the code
    frame=single, % Border style for the code block
    bgcolor=darkbackground, % Background color for the code block
    linenos, % Display line numbers
    numbersep=5pt, % Space between line numbers and code
    label=code-\thecodeblock}

% Custom Environments
\tcbuselibrary{listingsutf8}
\newtcolorbox[use counter=mydefinition, number within=section]{mydefinition}[1][]{colback=white, colframe=cyanaccent, fonttitle=\bfseries, breakable=true, enhanced, parbox=false, title=Definition~\themydefinition, #1}
\newtcolorbox[use counter=myidentity, number within=section]{myidentity}[1][]{colback=white, colframe=myorange, fonttitle=\bfseries, breakable=true, enhanced jigsaw, before skip=10pt, after skip=10pt, parbox=false, title=Identity~\themyidentity, #1}
\newtcolorbox[use counter=myproof, number within=section]{myproof}[1][]{colback=white, colframe=mygreen, fonttitle=\bfseries, breakable=true, enhanced jigsaw, before skip=10pt, after skip=10pt, parbox=false, title=Proof~\themyproof, #1}
\newtcolorbox[use counter=mynote, number within=section]{mynote}[1][]{colback=white, colframe=accent3, fonttitle=\bfseries, breakable=true, enhanced jigsaw, before skip=10pt, after skip=10pt, parbox=false, title=Note~\themynote, #1}
\newtcolorbox[use counter=mycodeblock, number within=section]{mycodeblock}[1][]{colback=darkbackground, colframe=accent2, fonttitle=\bfseries, breakable=unlimited, enhanced jigsaw, before skip=10pt, after skip=10pt, parbox=false, title=Code~\themycodeblock, #1}
\newtcolorbox[use counter=myexercise, number within=section]{myexercise}[1][]{colback=white, colframe=accent1, fonttitle=\bfseries, breakable=true, enhanced jigsaw, before skip=10pt, after skip=10pt, parbox=false, title=Exercise~\themyexercise, #1}

% Custom Theorem-like Environments
\newtheorem{mytheorem}{Theorem}[subsection]
\newtheorem{mydefinitiontheorem}{Definition}[subsection]
\newtheorem{myexercisetheorem}{Exercise}[subsection]

% Syntax Highlighting
\usemintedstyle{monokai}
\setminted{
    bgcolor=darkbackground,
    style=monokai,
    linenos,
    breaklines,
    fontsize=\small
}

% Custom Title Page with TikZ Decorations and Logo
\newcommand{\customtitle}[6]{
    \begin{titlepage}
        \noindent
        \begin{tikzpicture}[remember picture,overlay]
            \begin{scope}[blend mode=soft light]
                \fill[accent1] (current page.south west) rectangle (current page.north east);
                \fill[accent2] (0.2*\paperwidth,0.2*\paperheight) -- (current page.north east) -- (current page.north west) -- cycle;
                \fill[accent3] (current page.south west) -- (current page.north west) -- (0.8*\paperwidth,0.8*\paperheight) -- cycle;
                \draw[line width=1mm, cyanaccent, opacity=0.5] (current page.south west) -- (current page.north east);
                \draw[line width=1mm, cyanaccent, opacity=0.5] (current page.south east) -- (current page.north west);
            \end{scope}
            % Logo
            \node [anchor=north, yshift=-8.2cm] at (current page.north) {
                \includegraphics[width=0.5\textwidth]{#6}
            };
            % Title and Subtitle in cyan box
            \node [fill=cyanaccent, fill opacity=0.8, text opacity=1, inner sep=1cm, anchor=north, text width=\paperwidth, minimum height=0.15\paperheight, align=center] at ([yshift=-2cm]current page.center) {
                \Huge\textbf{\textit{#1}}\\[0.5cm]
                \Large\textbf{\textit{#5}}
            };
            % Author, Subject, and Date in royal blue box
            \node [fill=accent3, fill opacity=0.8, text opacity=1, inner sep=0.5cm, anchor=north, text width=\textwidth-8cm, align=center] at ([yshift=-6.45cm]current page.center) {
                \Large\textbf{\textit{#2}}\\[0.5cm]
                \normalsize\textbf{#3}\\[0.5cm]
                \normalsize\textbf{#4}
            };
        \end{tikzpicture}
    \end{titlepage}
}

% Custom Chapter Page with TikZ Decorations
\newcommand{\chapterpage}{
    \begin{tikzpicture}[remember picture,overlay]
        \begin{scope}[blend mode=soft light]
            \fill[accent1] (0,0) rectangle (\paperwidth,\paperheight);
            \fill[accent2] (0.2*\paperwidth,0.2*\paperheight) -- (\paperwidth,0.8*\paperheight) -- (\paperwidth,0) -- cycle;
            \fill[accent3] (0,0) -- (\paperwidth,0) -- (0.8*\paperwidth,0.8*\paperheight) -- cycle;
        \end{scope}
    \end{tikzpicture}
}

%Custom math commands
\newcommand{\intg}[3]{\int_{#1}^{#2} #3 \, dx}
\newcommand{\sumg}[3]{\sum_{#1}^{#2} #3}
\newcommand{\diff}[2]{\frac{d #1}{d #2}}
\newcommand{\pdiff}[2]{\frac{\partial #1}{\partial #2}} % Partial differential
\newcommand{\limit}[3]{\lim_{#1 \to #2} #3} % Limit
\newcommand{\infseries}[2]{\sum_{#1}^{\infty} #2} % Infinite series
\newcommand{\mat}[1]{\begin{matrix} #1 \end{matrix}} % Matrix
\newcommand{\pmat}[1]{\begin{pmatrix} #1 \end{pmatrix}} % Parentheses matrix
\newcommand{\bmat}[1]{\begin{bmatrix} #1 \end{bmatrix}} % Bracket matrix
\newcommand{\mycases}[1]{\left\{\begin{array}{lr}
        x(n), & \text{for } 0\leq n\leq 1\\
        x(n-1), & \text{for } 0\leq n\leq 1\\
        x(n-1), & \text{for } 0\leq n\leq 1
        \end{array}\right\}}
\newcommand{\floor}[1]{\left\lfloor #1 \right\rfloor} % Floor function
\newcommand{\ceil}[1]{\left\lceil #1 \right\rceil} % Ceiling function
\newcommand{\abs}[1]{\left| #1 \right|} % Absolute value
\newcommand{\norm}[1]{\left\| #1 \right\|} % Norm
\newcommand{\set}[1]{\left\{ #1 \right\}} % Set
\newcommand{\ang}[1]{\left\langle #1 \right\rangle} % Angle
\newcommand{\eqn}[1]{\begin{equation} #1 \end{equation}} % Equation
\newcommand{\eqnarrayx}[1]{\begin{eqnarray} #1 \end{eqnarray}} % Equation array
\newcommand{\R}{\mathbb{R}}
\newcommand{\Z}{\mathbb{Z}}
\newcommand{\N}{\mathbb{N}}
\newcommand{\Q}{\mathbb{Q}}
\newcommand{\forallx}{\forall x \in \mathbb{R}}
\newcommand{\existsx}{\exists x \in \mathbb{R}}
\newcommand{\iintd}[1]{\iint\!#1\,dx\,dy}
\newcommand{\iiintd}[1]{\iiint\!#1\,dx\,dy\,dz}
\newcommand{\nint}[3]{\int_{#1}^{#2}\!#3\,dx}

% Custom Headers and Footers
\pagestyle{fancy}
\fancyhf{}
\fancyhead[LE,RO]{\thepage}
\fancyhead[RE]{\leftmark}
\fancyhead[LO]{\rightmark}

% Section Formatting
\titleformat{\section}
  {\normalfont\Large\bfseries\color{cyanaccent}}
  {\thesection}{1em}{}
\titleformat{\subsection}
  {\normalfont\large\bfseries\color{cyanaccent}}
  {\thesubsection}{1em}{}
\titleformat{\subsubsection}
  {\normalfont\normalsize\bfseries\color{cyanaccent}}
  {\thesubsubsection}{1em}{}

\newcommand{\settitle}[1]{
\hypersetup{
    colorlinks=true,
    linkcolor=cyanaccent,
    filecolor=cyanaccent,
    urlcolor=cyanaccent,
    citecolor=cyanaccent,
    pdftitle={#1},
    pdfpagemode=FullScreen
}}

\settitle{Demonstration Document}
\begin{document}

% Custom Cover Page
\customtitle{Latex Template}{A demonstration}{Eklavya Raman}{\today}{Test}{sample.jpg}

% Table of Contents
\tableofcontents
\newpage

% Chapter with integrated chapter page
\chapter{Introduction}
\chapterpage
\section{Background}
\lipsum[1]
\subsection{History}
\lipsum[2]

% Another Chapter
\chapter{Main Content}
\chapterpage
\section{Definitions}
\begin{mydefinition}
\lipsum[1-6]
\end{mydefinition}

\section{Identities}
\begin{myidentity}
This is a custom identity box.
\begin{enumerate}
    \item Here is a floor function: $\floor{x}$
    \item Here is a ceiling function: $\ceil{x}$
    \item Here is an integral: $\intg{a}{b}{f(x)}$
    \item Here is a sum: $\sumg{n=1}{\infty}{a_n}$
    \item Here is a differential: $\diff{y}{x}$
    \item Here is a partial differential: $\pdiff{u}{x}$
    \item Here is a limit: $\limit{n}{\infty}{a_n}$
    \item Here is an infinite series: $\infseries{n=1}{a_n}$
    \item Here is a matrix: $\mat{a & b \\ c & d}$
    \item Here is a parentheses matrix: $\pmat{a & b \\ c & d}$
    \item Here is a bracket matrix: $\bmat{a & b \\ c & d}$
    \item Here are some cases: 
\[
\mycases{
  1 & \text{if} x \ge 0 \\
  -1 & \text{if} x < 0
}
\]
    \item Here is an absolute value: $\abs{x}$
    \item Here is a norm: $\norm{x}$
    \item Here is a set: $\set{x \mid x \ge 0}$
    \item Here is an angle: $\ang{x}$
    \item Here is an equation: \eqn{E = mc^2}
    \item Here is an equation array: 
\eqnarrayx{
  a &=& b + c \\
  d &=& e + f
  }
    \newpage
    \item Here is a double integral: $\iintd{3x+4y}$
    \item Here is a triple integral: $\iiintd{3x+4y+5z}$
    \item Here are some more commands: $\R, \Q, \Z, \N, \forallx, \existsx$
\end{enumerate}
\end{myidentity}

\section{Proofs}
\begin{myproof}
This is a custom proof box.
\end{myproof}

\section{Notes}
\begin{mynote}
This is a custom note box.
\end{mynote}

\section{Code Blocks}
\subsection{Custom Code block}
\begin{longlisting}
\begin{minted}{python}
import math
import numpy
from PIL import Image


def ascii_color_encode():
    ascii_colors = {}

    # Generating unique RGB values for ASCII characters
    for i in range(128):
        red = (i * 17) % 256  # Unique red component
        green = (i * 37) % 256  # Unique green component
        blue = (i * 23) % 256  # Unique blue component

        # Storing RGB values in a dictionary with ASCII value as key
        ascii_colors[i] = [red, green, blue]

    # Displaying the dictionary
    return ascii_colors


def text_to_picture(path, ascii_colours, save_as):
    file = open(path, 'r')
    data = file.read()
    ascii_values = []
    for i in data:
        if ord(i) < 128:  # Check if character falls within ASCII range
            ascii_values.append(ord(i))
        else:
            ascii_values.append(ord(' '))  # Replace non-ASCII characters with a space ' '
    square = int(math.ceil(math.sqrt(len(ascii_values))))
    print(square)  # Adjusted square size
    total_elements = square * square
    for i in range(total_elements - len(ascii_values)):
        ascii_values.append(ord(' '))  # Fill remaining spaces with ' '
    pixel_array = []
    for i in range(0, len(ascii_values)):
        pixel_array.append(ascii_colours[ascii_values[i]])

    array = numpy.array(pixel_array, dtype=numpy.uint8).reshape(square, square, 3)

    image = Image.fromarray(array)
    image.resize((1000, 1000))
    image.save(f'{save_as}.png')


def picture_to_text(picture_path, ascii_colours, output_name):
    image = Image.open(picture_path)
    pixels = image.load()
    width, height = image.size
    data = []
    for i in range(width):
        for j in range(height):
            data.append(pixels[j, i])
    text = ''
    for k in data:
        for j in ascii_colours:
            if ascii_colours[j] == list(k):
                text += chr(j)

    with open(f'{output_name}.txt', 'w') as file:
        file.write(text)
text_to_picture(path='sample.txt', ascii_colours=ascii_color_encode(), save_as='image')}
\end{minted}
\end{longlisting}
\subsection{Source code for template}
\begin{longlisting}
\begin{minted}{latex}
\usepackage[a4paper, margin=1in]{geometry}
\usepackage[margin=0pt, noheadfoot]{geometry}
\usepackage{minted}
\setminted{breaklines, % Break lines automatically
    breakanywhere, % Allow breaks anywhere
    fontsize=\small, % Font size for the code
    frame=single, % Border style for the code block
    bgcolor=darkbackground, % Background color for the code block
    linenos, % Display line numbers
    numbersep=5pt, % Space between line numbers and code
    highlightlines=1,% Highlight specific lines (if any)
    label=code-\thecodeblock}
\usepackage{caption}
\usepackage{tcolorbox}
\tcbuselibrary{breakable, skins, theorems, listings}
\usepackage{titlesec}
\usepackage{fancyhdr}
\usepackage{hyperref}
\usepackage{lipsum}
\usepackage{amsmath}
\usepackage{amsthm}
\usepackage{amssymb}
\usepackage{enumitem}
\usepackage{tikz}
\usepackage{tocloft}

% Colors
\usepackage{xcolor}
\definecolor{darkbackground}{HTML}{121212}
\definecolor{whitetext}{HTML}{FFFFFF}
\definecolor{cyanaccent}{HTML}{00FFFF}
\definecolor{lightcyan}{HTML}{E0FFFF}
\definecolor{lightgray}{HTML}{D3D3D3}
\definecolor{accent1}{HTML}{0047AB}
\definecolor{accent2}{HTML}{6495ED}
\definecolor{accent3}{HTML}{3F00FF}
\definecolor{myorange}{HTML}{FFA500}
\definecolor{mylightblue}{HTML}{ADD8E6}
\definecolor{mygreen}{HTML}{90EE90}

\pagecolor{darkbackground}
\color{whitetext}

\newenvironment{longlisting}{\captionsetup{type=listing}}{}

% Custom Counters
\newcounter{mydefinition}[subsection]
\renewcommand{\themydefinition}{\thesubsection.\arabic{mydefinition}}
\newcounter{myidentity}[subsection]
\renewcommand{\themyidentity}{\thesubsection.\arabic{myidentity}}
\newcounter{myproof}[subsection]
\renewcommand{\themyproof}{\thesubsection.\arabic{myproof}}
\newcounter{mynote}[subsection]
\renewcommand{\themynote}{\thesubsection.\arabic{mynote}}
\newcounter{mycodeblock}[subsection]
\renewcommand{\themycodeblock}{\thesubsection.\arabic{mycodeblock}}
\newcounter{myexercise}[subsection]
\renewcommand{\themyexercise}{\thesubsection.\arabic{myexercise}}
\newcounter{codeblock}[section]
\renewcommand{\thecodeblock}{\thesubsection.\arabic{codeblock}}


\setminted{breaklines, % Break lines automatically
    breakanywhere, % Allow breaks anywhere
    fontsize=\small, % Font size for the code
    frame=single, % Border style for the code block
    bgcolor=darkbackground, % Background color for the code block
    linenos, % Display line numbers
    numbersep=5pt, % Space between line numbers and code
    label=code-\thecodeblock}

% Custom Environments
\tcbuselibrary{listingsutf8}
\newtcolorbox[use counter=mydefinition, number within=section]{mydefinition}[1][]{colback=white, colframe=cyanaccent, fonttitle=\bfseries, breakable=true, enhanced, parbox=false, title=Definition~\themydefinition, #1}
\newtcolorbox[use counter=myidentity, number within=section]{myidentity}[1][]{colback=white, colframe=myorange, fonttitle=\bfseries, breakable=true, enhanced jigsaw, before skip=10pt, after skip=10pt, parbox=false, title=Identity~\themyidentity, #1}
\newtcolorbox[use counter=myproof, number within=section]{myproof}[1][]{colback=white, colframe=mygreen, fonttitle=\bfseries, breakable=true, enhanced jigsaw, before skip=10pt, after skip=10pt, parbox=false, title=Proof~\themyproof, #1}
\newtcolorbox[use counter=mynote, number within=section]{mynote}[1][]{colback=white, colframe=accent3, fonttitle=\bfseries, breakable=true, enhanced jigsaw, before skip=10pt, after skip=10pt, parbox=false, title=Note~\themynote, #1}
\newtcolorbox[use counter=mycodeblock, number within=section]{mycodeblock}[1][]{colback=darkbackground, colframe=accent2, fonttitle=\bfseries, breakable=unlimited, enhanced jigsaw, before skip=10pt, after skip=10pt, parbox=false, title=Code~\themycodeblock, #1}
\newtcolorbox[use counter=myexercise, number within=section]{myexercise}[1][]{colback=white, colframe=accent1, fonttitle=\bfseries, breakable=true, enhanced jigsaw, before skip=10pt, after skip=10pt, parbox=false, title=Exercise~\themyexercise, #1}

% Custom Theorem-like Environments
\newtheorem{mytheorem}{Theorem}[subsection]
\newtheorem{mydefinitiontheorem}{Definition}[subsection]
\newtheorem{myexercisetheorem}{Exercise}[subsection]

% Syntax Highlighting
\usemintedstyle{monokai}
\setminted{
    bgcolor=darkbackground,
    style=monokai,
    linenos,
    breaklines,
    fontsize=\small
}

% Custom Title Page with TikZ Decorations and Logo
\newcommand{\customtitle}[6]{
    \begin{titlepage}
        \noindent
        \begin{tikzpicture}[remember picture,overlay]
            \begin{scope}[blend mode=soft light]
                \fill[accent1] (current page.south west) rectangle (current page.north east);
                \fill[accent2] (0.2*\paperwidth,0.2*\paperheight) -- (current page.north east) -- (current page.north west) -- cycle;
                \fill[accent3] (current page.south west) -- (current page.north west) -- (0.8*\paperwidth,0.8*\paperheight) -- cycle;
                \draw[line width=1mm, cyanaccent, opacity=0.5] (current page.south west) -- (current page.north east);
                \draw[line width=1mm, cyanaccent, opacity=0.5] (current page.south east) -- (current page.north west);
            \end{scope}
            % Logo
            \node [anchor=north, yshift=-8.2cm] at (current page.north) {
                \includegraphics[width=0.5\textwidth]{#6}
            };
            % Title and Subtitle in cyan box
            \node [fill=cyanaccent, fill opacity=0.8, text opacity=1, inner sep=1cm, anchor=north, text width=\paperwidth, minimum height=0.15\paperheight, align=center] at ([yshift=-2cm]current page.center) {
                \Huge\textbf{\textit{#1}}\\[0.5cm]
                \Large\textbf{\textit{#5}}
            };
            % Author, Subject, and Date in royal blue box
            \node [fill=accent3, fill opacity=0.8, text opacity=1, inner sep=0.5cm, anchor=north, text width=\textwidth-8cm, align=center] at ([yshift=-6.45cm]current page.center) {
                \Large\textbf{\textit{#2}}\\[0.5cm]
                \normalsize\textbf{#3}\\[0.5cm]
                \normalsize\textbf{#4}
            };
        \end{tikzpicture}
    \end{titlepage}
}

% Custom Chapter Page with TikZ Decorations
\newcommand{\chapterpage}{
    \begin{tikzpicture}[remember picture,overlay]
        \begin{scope}[blend mode=soft light]
            \fill[accent1] (0,0) rectangle (\paperwidth,\paperheight);
            \fill[accent2] (0.2*\paperwidth,0.2*\paperheight) -- (\paperwidth,0.8*\paperheight) -- (\paperwidth,0) -- cycle;
            \fill[accent3] (0,0) -- (\paperwidth,0) -- (0.8*\paperwidth,0.8*\paperheight) -- cycle;
        \end{scope}
    \end{tikzpicture}
}

%Custom math commands
\newcommand{\intg}[3]{\int_{#1}^{#2} #3 \, dx}
\newcommand{\sumg}[3]{\sum_{#1}^{#2} #3}
\newcommand{\diff}[2]{\frac{d #1}{d #2}}
\newcommand{\pdiff}[2]{\frac{\partial #1}{\partial #2}} % Partial differential
\newcommand{\limit}[3]{\lim_{#1 \to #2} #3} % Limit
\newcommand{\infseries}[2]{\sum_{#1}^{\infty} #2} % Infinite series
\newcommand{\mat}[1]{\begin{matrix} #1 \end{matrix}} % Matrix
\newcommand{\pmat}[1]{\begin{pmatrix} #1 \end{pmatrix}} % Parentheses matrix
\newcommand{\bmat}[1]{\begin{bmatrix} #1 \end{bmatrix}} % Bracket matrix
\newcommand{\mycases}[1]{\left\{\begin{array}{lr}
        x(n), & \text{for } 0\leq n\leq 1\\
        x(n-1), & \text{for } 0\leq n\leq 1\\
        x(n-1), & \text{for } 0\leq n\leq 1
        \end{array}\right\}}
\newcommand{\floor}[1]{\left\lfloor #1 \right\rfloor} % Floor function
\newcommand{\ceil}[1]{\left\lceil #1 \right\rceil} % Ceiling function
\newcommand{\abs}[1]{\left| #1 \right|} % Absolute value
\newcommand{\norm}[1]{\left\| #1 \right\|} % Norm
\newcommand{\set}[1]{\left\{ #1 \right\}} % Set
\newcommand{\ang}[1]{\left\langle #1 \right\rangle} % Angle
\newcommand{\eqn}[1]{\begin{equation} #1 \end{equation}} % Equation
\newcommand{\eqnarrayx}[1]{\begin{eqnarray} #1 \end{eqnarray}} % Equation array
\newcommand{\R}{\mathbb{R}}
\newcommand{\Z}{\mathbb{Z}}
\newcommand{\N}{\mathbb{N}}
\newcommand{\Q}{\mathbb{Q}}
\newcommand{\forallx}{\forall x \in \mathbb{R}}
\newcommand{\existsx}{\exists x \in \mathbb{R}}
\newcommand{\iintd}[1]{\iint\!#1\,dx\,dy}
\newcommand{\iiintd}[1]{\iiint\!#1\,dx\,dy\,dz}
\newcommand{\nint}[3]{\int_{#1}^{#2}\!#3\,dx}

% Custom Headers and Footers
\pagestyle{fancy}
\fancyhf{}
\fancyhead[LE,RO]{\thepage}
\fancyhead[RE]{\leftmark}
\fancyhead[LO]{\rightmark}

% Section Formatting
\titleformat{\section}
  {\normalfont\Large\bfseries\color{cyanaccent}}
  {\thesection}{1em}{}
\titleformat{\subsection}
  {\normalfont\large\bfseries\color{cyanaccent}}
  {\thesubsection}{1em}{}
\titleformat{\subsubsection}
  {\normalfont\normalsize\bfseries\color{cyanaccent}}
  {\thesubsubsection}{1em}{}

\newcommand{\settitle}[1]{
\hypersetup{
    colorlinks=true,
    linkcolor=cyanaccent,
    filecolor=cyanaccent,
    urlcolor=cyanaccent,
    citecolor=cyanaccent,
    pdftitle={#1},
    pdfpagemode=FullScreen
}}
\end{minted}
\end{longlisting}

\section{Exercises}
\begin{myexercise}
This is a custom exercise box.
\end{myexercise}

\begin{myexercise}
Another exercise box.
\end{myexercise}

\chapter{Test}
\chapterpage
\lipsum[3-4]


\end{document}
